\documentclass[12pt]{article}
\usepackage[utf8]{inputenc}
\usepackage[T1]{fontenc}
\usepackage{hyperref}
\title{Talinga-Bwisi Language Resources for K-12 Curriculum Development}
\author{Compiled by Claude Sonnet 5 (claude-sonnet-5) for Andr\'as Kornai}
\date{}
\begin{document}
\maketitle

\section{Introduction}
Talinga-Bwisi (ISO 639-3 tlj) is a Bantu language of the Uganda--DR Congo border region. It is called Talinga or Kitalinga in the DRC and Bwisi, Lubwisi or Olubwisi in Uganda. Its main areas are Bundibugyo District in Uganda and North Kivu in the DRC. Speaker estimates vary by source: about 99,000 (Wikipedia), and about 68,500 in Uganda plus 30,000 in the DRC (Omniglot). Documentation is described as very scant. This search found no school curriculum materials.

\section{Primary Educational Resources}
No primers or textbooks turned up in this search. The Lubwisi Spelling Guide (Kwamba/Lubwisi Development Association, third edition) is mentioned in the search summary. I did not see a URL for it, so I do not list it as a resource.

\section{Linguistic and Reference Materials}
\begin{itemize}
\item Wikipedia overview of the Talinga language: \url{https://en.wikipedia.org/wiki/Talinga_language}
\item Omniglot page on Talinga (Kitalinga / Lubwisi), with writing system information: \url{https://www.omniglot.com/writing/talinga.htm}
\end{itemize}

\section{Religious and Community Resources}
\begin{itemize}
\item Scripture Earth page for tlj. Per the search summary it lists a Gospel of Mark, a JESUS Film version, a Bible app and audio recordings: \url{https://www.scriptureearth.org/00eng.php?iso=tlj}
\end{itemize}

\section{Digital and Audio Resources}
The Scripture Earth page above links to Global Recordings Network audio, according to the search summary. I did not see the audio URLs themselves.

\section{Curriculum Development Context}
Very little exists beyond scripture, a spelling guide and brief reference pages. A curriculum would need new material, built with the community and the Kwamba/Lubwisi Development Association. The Lubwisi orthography is Latin-based and was developed in the 1990s.

\end{document}